\clearpage
\section{Hilfsmittelverzeichnis}
\label{sec:hilfsmittel}

\begin{table}[htpb]
    \centering
    \caption{Auflistung der verwendeten Hilfsmittel}
    \label{tab:hilfsmittelverzeichnis}
    \begin{tabularx}{\textwidth}{>{\raggedright\arraybackslash}X
                                 >{\raggedright\arraybackslash}X
                                 >{\raggedright\arraybackslash}X
                                 l}
        \toprule
        \textbf{Hilfsmittel} & \textbf{Einsatzform} & \textbf{Betroffene Teile der Arbeit} & \textbf{Datum} \\
        \midrule
        Google (Scholar) (keine KI) & Gezielte Suche nach Papers zu ZTD & Kapitel 2 & Oktober 2026 \\
        \rowcolor{lightgray}
        Sonar (Perplexity) & Einmalige Suche von DOI von Publikationen & Kapitel 2 & Oktober 2026 \\
        Claude Sonnet 5.5 (Anthropic) & Umformulierung, Anpassung und Überprüfung durch Autor*in & Kapitel 2.4 & Oktober 2026 \\
        \rowcolor{lightgray}
        Claude Sonnet 5.5 (Anthropic) & Formatierung und Erstellung von \LaTeX{} für Tabellen & Kapitel 3 und 8 & Sept.~--~Nov. 2026 \\
        Claude Opus 5.5 (Anthropic) & Coding zum Aufsetzen des Monitorings und zur Zeitreihenanalyse & Kapitel 5 & November 2026 \\
        \rowcolor{lightgray}
        Claude Sonnet 5.5 (Anthropic) & Syntax- und Grammatikprüfung & Ganzer Bericht & November 2026 \\
        \bottomrule
    \end{tabularx}
\end{table}